% Generated draft: requires mathematical and source review.
\documentclass[11pt]{article}
\usepackage[margin=25mm]{geometry}
\usepackage{fontspec}
\setmainfont{Latin Modern Roman}
\usepackage{amsmath,amssymb,mathtools,unicode-math}
\setmathfont{Latin Modern Math}
\usepackage{enumitem,needspace,xurl,hyperref}
\setlength{\parindent}{0pt}
\setlength{\parskip}{5pt}
\setlength{\emergencystretch}{3em}
\newcommand{\N}{\mathbb{N}}
\newcommand{\Z}{\mathbb{Z}}
\newcommand{\Q}{\mathbb{Q}}
\newcommand{\R}{\mathbb{R}}
\newcommand{\C}{\mathbb{C}}
\newcommand{\F}{\mathbb{F}}
\newcommand{\A}{\mathbb{A}}
\newcommand{\PP}{\mathbb{P}}
\newcommand{\E}{\mathbb{E}}
\newcommand{\eps}{\varepsilon}
\newcommand{\dd}{\,\mathrm{d}}
\begin{document}
\section*{UnsolvedMath ID: 30002698}
\textbf{Source problem code:} catalogue entry
\subsection{Definitions and mathematical statement}
For a smooth compact manifold \(M\), write \(b_2(M)=\dim_{\R}H^2(M;\R)\). In the usual compact homogeneous convention, a generalized flag manifold is a quotient \(G/K\), where \(G\) is a compact connected semisimple Lie group and \(K\) is the centralizer of a torus in \(G\). The supplied record does not specify its precise group-theoretic convention; this definition explains the terminology without resolving that omission. A Kähler--Einstein metric is a Kähler metric \(g\) satisfying \(\operatorname{Ric}(g)=\lambda g\) for a real constant \(\lambda\).

Consider a subgroup chain \(K\subset L\subset G\) giving the homogeneous fibration
\[
L/K\longrightarrow G/K\xrightarrow{\ \pi\ }G/L,
\qquad \pi(gK)=gL,
\]
where both \(G/K\) and \(G/L\) are generalized flag manifolds. The isotropy representation of \(G/K\) is the action of \(K\) on \(T_{eK}(G/K)\); an isotropy decomposition expresses this tangent space as a direct sum of irreducible isotropy modules.

The supplied statement reports that a method using a Kähler--Einstein metric and a homogeneous-fibration Ricci-tensor formula determines the Ricci tensor in the setting
\[
b_2(G/K)=2,\qquad b_2(G/L)=1,
\]
with at most six isotropy summands. The record does not explicitly identify which space's isotropy decomposition is being counted. This reported scope is supplied background, not a theorem established by this draft.

\textbf{Question.} Does the fibration method using Kähler--Einstein metrics and the Ricci-tensor formula referred to as ``Lemma 2'' in the original statement work for every generalized flag manifold, extending beyond the stated restrictions?

The universal target is every generalized flag manifold. However, the record does not reproduce Lemma 2, specify the family of metrics whose Ricci tensors are to be determined, or define what successful application of the method requires. It also does not say whether a single suitable fibration, a sequence of fibrations, or additional data are permitted. Consequently, the question can be preserved as a methodological extension problem, but a fully specified quantified conjecture cannot be extracted from the supplied material alone.
\subsection{Short English statement}
The problem asks whether a method using homogeneous fibrations and Kähler--Einstein metrics can determine Ricci tensors for every generalized flag manifold. The supplied special setting has second Betti numbers two and one and at most six isotropy summands. The missing Lemma 2 and unspecified meaning of successful extension require human review.
\subsection{Sources}
\begin{itemize}
\item Local UnsolvedMath registry: supplied complete record for permanent ID \(30002698\), entitled \emph{Ricci-Tensor Fibrations for Generalized Flag Manifolds}; exact source problem code \texttt{catalogue entry}. The record gives proposed year \(2014\) and names no proposer.
\item Supplied source citation: \emph{Mini-Workshop: Einstein Metrics, Ricci Solitons and Ricci Flow under Symmetry Assumptions (2015)}, Oberwolfach Reports, pp.~2529--2568, DOI \url{https://doi.org/10.4171/owr/2014/45}. The literature list identifies this as \emph{Oberwolfach Report 45/2014}. The supplied year discrepancy remains unresolved.
\item Supplied literature source: \emph{Flag Manifolds with Five Isotropy Summands}, \url{https://doi.org/10.48550/arXiv.1207.2897}.
\item Supplied literature source: \emph{Classification of Einstein Metrics on Flag Manifolds}, \url{https://doi.org/10.1016/j.bulsci.2013.11.002}.
\item Literature assessment dated exactly \texttt{2026-08-21}: the supplied curation labels the problem partially solved and states that the fibration and Riemannian-submersion method has classified numerous flag-manifold families, including five-isotropy-summand cases, but has not provided a complete method or classification for every generalized flag manifold. This is supplied dated curation derived from the Oberwolfach material, not an independently verified literature assessment; its status and cited claims require independent checking.
\item The supplied record labels its statement \texttt{corrected\_verified} and reports comparison with the official Oberwolfach report. That verification claim is source metadata and has not been independently checked in preparing this draft.
\end{itemize}
\end{document}
